% Chapter 6: Model checking
% Bayesian Data Analysis - Gelman et al.

\chapter{模型检查}

\section{本章导论}

在前几章中，我们学习了贝叶斯推断的核心工具——如何从先验和似然出发，计算后验分布，以及如何构建层次模型来合并来自不同来源的信息。然而，在实际应用中，一个关键问题始终萦绕着我们：\textbf{我们如何知道所选的模型是否合适？}

模型检查是贝叶斯数据分析中一个容易被忽视但至关重要的环节。在经典统计中，我们有各种拟合优度检验；在贝叶斯框架中，我们同样需要工具来评估模型与数据的契合程度。本章将介绍一种强有力的方法——\textbf{后验预测检查（posterior predictive checking）}，它允许我们在贝叶斯框架内系统地评估模型。

我们将看到，模型检查不仅是验证模型是否合理，更是理解模型局限性、发现改进方向的重要途径。一个好的模型检查可以揭示数据中与模型假设不符的模式，从而引导我们改进模型或收集更多数据。

\section{后验预测检查的思想}

\subsection{为什么需要模型检查？}

贝叶斯分析的最终目标不是找到一个``正确''的模型，而是找到一个\textbf{有用}的模型——一个能够捕捉数据关键特征、产生合理预测的模型。任何模型都是对现实的简化，因此模型检查的核心问题是：\textbf{模型与数据的契合程度是否足以满足我们的分析目的？}

考虑一个简单例子。假设我们拟合了一个正态模型 $y_i \sim \mathrm{N}(\mu, \sigma^2)$ 到一组数据，但我们如何知道正态性假设是否合理？或者，当我们使用层次模型时，我们假设组内参数是可交换的——但数据是否支持这个假设？

这些问题的答案不能仅从模型内部获得，而需要我们主动将模型预测与实际数据进行比较。

\subsection{后验预测分布}

后验预测分布是模型检查的基础。给定观测数据 $y$，未来或新数据 $\tilde{y}$ 的后验预测分布为：

\[
p(\tilde{y} | y) = \int p(\tilde{y} | \theta) p(\theta | y) \, d\theta
\]

这个分布整合了参数的不确定性，反映了给定当前数据后模型对未来的预测。

\textbf{关键思想}是：如果模型合理，那么\textbf{真实数据应该与后验预测分布一致}。换句话说，我们不应该期望观测数据是后验预测分布的异常值。

\subsection{后验预测检查的直觉}

后验预测检查的基本思路如下：

\begin{enumerate}
  \item 从后验分布 $p(\theta | y)$ 中抽取参数值 $\theta^s$
  \item 给定每个 $\theta^s$，从似然模型 $p(\tilde{y} | \theta^s)$ 抽取伪数据 $\tilde{y}^s$
  \item 比较真实的观测数据 $y$ 与后验预测分布 $p(\tilde{y} | y)$
\end{enumerate}

如果 $y$ 落在后验预测分布的中心区域，模型与数据一致；如果 $y$ 落在尾部，则模型可能存在问题。

\begin{Example}[害虫治理研究]
  这是一个来自农业研究的例子。研究者比较了八种害虫栖息地的治理效果。数据包括每栖息地昆虫数量的计数。

  假设我们拟合了一个 Poisson 模型：$y_i \sim \mathrm{Poisson}(\theta_i)$，其中 $\theta_i$ 是栖息地 $i$ 的期望昆虫数量。

  通过后验预测检查，我们可以：
  \begin{itemize}
    \item 对每个后验样本 $\theta^s$，生成伪数据 $\tilde{y}^s$
    \item 比较观测数据 $y$ 与伪数据分布
    \item 检查是否存在异常大的观测值（overdispersed data）
  \end{itemize}

  如果某些观测值在后验预测分布的尾部，这可能表明存在过度离散（overdispersion），提示我们应该使用负二项分布或其他模型。
\end{Example}

\section{验证统计量与模型检查量}

\subsection{验证统计量}

验证统计量（test quantity）或模型检查量（discrepancy measure）$T(y, \theta)$ 是将数据和参数映射到标量或向量的函数，用于衡量模型与数据的某个特定方面。

常见的验证统计量包括：

\begin{itemize}
  \item \textbf{样本均值}：$T(y, \theta) = \bar{y}$
  \item \textbf{样本方差}：$T(y, \theta) = s^2 = \frac{1}{n-1}\sum_{i=1}^n (y_i - \bar{y})^2$
  \item \textbf{最大值}：$T(y, \theta) = \max(y_i)$
  \item \textbf{相关系数}：$T(y, \theta) = \text{corr}(y_i, y_j)$
\end{itemize}

在模型检查中，我们比较 $T(y, \theta)$ 与其后验预测分布 $p(T(\tilde{y}, \theta) | y)$。

\subsection{后验预测 p 值}

后验预测 p 值是后验预测检查的正式化表述。对于验证统计量 $T$，定义：

\[
\text{p-value} = \Pr(T(\tilde{y}, \theta) \geq T(y, \theta) | y)
\]

更精确地说，在参数空间上计算：

\[
\text{p-value} = \int \mathbb{I}(T(\tilde{y}, \theta^s) \geq T(y, \theta^s)) p(\theta^s | y) \, d\theta^s
\]

其中 $\theta^s$ 是从后验分布抽取的样本。

\textbf{解释}：后验预测 p 值度量的是，在给定数据和模型下，伪数据产生比观测数据更极端验证统计量的概率。如果这个概率很小（例如小于 0.05），则表明模型与数据在 $T$ 这个方面不一致。

\begin{Remark}
  后验预测 p 值与经典 p 值有本质区别。经典 p 值假设数据是随机化的，而\textbf{后验预测 p 值是条件于观测数据的}，反映的是模型对观测数据的拟合程度。
\end{Remark}

\section{模型检查的实践方法}

\subsection{检查批量}

当模型包含多个组或多个观测时，我们经常需要检查各个批量（batch）。设 $y_j$ 是第 $j$ 个组的观测，$\theta_j$ 是对应参数。

\textbf{批量特异性检查}可以揭示模型在哪些方面存在问题：

\begin{enumerate}
  \item \textbf{参数批量检查}：对每个组 $j$，计算 $T(\theta_j, \theta_{-j})$，其中 $\theta_{-j}$ 表示除 $j$ 以外的所有参数
  \item \textbf{数据批量检查}：对每个组 $j$，计算验证统计量 $T(y_j, \theta)$，然后比较 $T(y_j, \theta)$ 与其后验预测分布
\end{enumerate}

\begin{Example}[多组数据的批量检查]
  考虑一个分层模型，$J$ 所学校的学生成绩。每所学校有自己的真实效应 $\theta_j$，数据 $y_j$ 是带有测量误差的估计。

  批量检查可以帮助我们识别：
  \begin{itemize}
    \item 哪些学校的估计与模型预测差异最大
    \item 哪些学校的方差与层次模型假设不符
    \item 是否存在异常值学校
  \end{itemize}

  通过这种方法，我们不仅检查整体模型拟合，还可以定位具体问题所在。
\end{Example}

\subsection{检查交互作用}

当模型包含变量之间的交互作用时，检查交互作用是否与数据一致非常重要。

设我们有二元变量 $z_i$ 和 $w_i$，以及它们的交互作用 $z_i w_i$。一个常见的检查是验证交互作用是否在预测 $y_i$ 中起到预期作用。

具体来说，我们可以：
\begin{enumerate}
  \item 在没有交互作用的模型下计算后验预测分布
  \item 检查包含交互作用是否显著改善预测
  \item 或者直接检查数据中交互作用模式是否与模型一致
\end{enumerate}

\section{过度离散与计数数据的模型检查}

\subsection{过度离散问题}

计数数据经常表现出比 Poisson 分布预测的更大的变异性，这称为\textbf{过度离散（overdispersion）}。

在 Poisson 模型中，方差等于均值：$\text{var}(y) = \mathbb{E}(y)$。但实际数据通常表现出更大的方差。

\subsection{检测过度离散}

后验预测检查可以有效检测过度离散：

\begin{enumerate}
  \item 计算观测数据的方差：$T(y, \theta) = \text{var}(y)$
  \item 从后验预测分布生成伪数据 $\tilde{y}^s$
  \item 计算每个 $\tilde{y}^s$ 的方差：$T(\tilde{y}^s, \theta^s)$
  \item 比较观测方差与后验预测方差分布
\end{enumerate}

如果观测方差在伪数据方差的尾部，则表明存在过度离散。

\begin{Example}[害虫数据的过度离散检查]
  回到害虫治理研究。Poisson 模型假设方差等于均值。但如果数据存在过度离散，观测方差将显著大于后验预测方差。

  通过后验预测检查，我们可以量化这种不一致程度。如果后验预测 p 值很小，我们可以考虑使用负二项分布或零膨胀模型。
\end{Example}

\section{外部数据验证}

后验预测检查的一种重要扩展是使用外部数据验证。核心思想是：如果模型是合理的，它应该能够预测我们尚未看到的新数据。

设我们有两个相关的数据集 $y_{\text{old}}$ 和 $y_{\text{new}}$。我们可以：

\begin{enumerate}
  \item 使用 $y_{\text{old}}$ 拟合模型，获得后验分布 $p(\theta | y_{\text{old}})$
  \item 使用后验预测分布预测 $y_{\text{new}}$：$p(y_{\text{new}} | y_{\text{old}}) = \int p(y_{\text{new}} | \theta) p(\theta | y_{\text{old}}) \, d\theta$
  \item 比较预测分布与实际观测的 $y_{\text{new}}$
\end{enumerate}

这种方法在时间序列预测、跨研究汇总等场景中特别有用。

\textbf{先验预测检查}是另一种相关方法，它在收集数据之前使用先验分布预测数据。这有助于：

\begin{itemize}
  \item 检查先验分布的合理性
  \item 理解模型对数据的预期
  \item 发现模型中的潜在问题
\end{itemize}

\section{模型检查与其他推断的关系}

\subsection{与模型选择的联系}

模型检查是模型选择的前置步骤。在比较多个模型时，我们不仅应该比较它们的预测性能，还应该检查每个模型与数据的兼容性。

一个在训练集上表现良好但在检查中暴露问题的模型，可能存在过拟合或遗漏关键特征。

\subsection{与敏感性分析的联系}

模型检查与敏感性分析有密切关系。当检查发现模型与数据不一致时，敏感性分析可以帮助我们理解这种不一致是否会影响我们的结论。

例如，如果后验预测检查表明存在过度离散，但我们关心的参数估计在不同的过度离散模型下保持稳定，那么模型的不完美可能不会影响我们的最终结论。

\section{模型检查的局限性}

虽然后验预测检查是一种强有力的工具，但它也有局限性：

\begin{enumerate}
  \item \textbf{检查量的选择}：我们只能检查我们想到的验证统计量。可能存在模型缺陷，但我们的检查量没有捕捉到。
  \item \textbf{多重比较问题}：如果我们检查很多验证统计量，一些小的 p 值可能是偶然的。
  \item \textbf{模型不确定性}：后验预测检查没有考虑模型本身的不确定性。
  \item \textbf{先验敏感性}：检查结果可能对先验选择敏感。
\end{enumerate}

\begin{Remark}
  这些局限性并不意味着我们应该放弃模型检查，而是提醒我们：模型检查应该被视为一种\textbf{探索性工具}，而非最终裁决。它帮助我们理解模型的 strengths 和 weaknesses，但最终的数据分析和决策需要综合考虑多方面因素。
\end{Remark}

\section{从频率派角度理解后验预测检查}

从频率派的视角，后验预测检查可以被理解为一种校准（calibration）方法。

如果模型是正确指定的，那么：
\begin{itemize}
  \item 对于每个验证统计量 $T$，后验预测 p 值在 $[0, 1]$ 上均匀分布
  \item 重复实验将产生均匀分布的 p 值
\end{itemize}

这种校准性质意味着，如果我们在许多不同数据集上进行后验预测检查，p 值应该均匀分布。如果 p 值系统性地偏向极端值（例如过多的小 p 值），则表明模型存在系统性问题。

从频率派视角看，后验预测检查可以被视为评估模型\textbf{校准性}的一种方法——模型是否正确地量化了不确定性。

\section{本章小结}

本章我们学习了模型检查的核心概念和方法：

\begin{enumerate}
  \item \textbf{模型检查的必要性}：任何模型都是对现实的简化，我们需要工具来评估模型与数据的契合程度
  \item \textbf{后验预测检查的思想}：通过比较观测数据与模型对数据的预测来评估模型
  \item \textbf{验证统计量与模型检查量}：用于衡量模型特定方面的函数
  \item \textbf{后验预测 p 值}：量化模型与数据不一致程度的正式指标
  \item \textbf{批量检查}：检查模型在不同数据子集上的表现
  \item \textbf{过度离散检测}：使用后验预测方法检测计数数据的过度变异
  \item \textbf{外部数据验证}：使用独立数据验证模型预测能力
  \item \textbf{模型检查的局限性}：检查量选择、多重比较、模型不确定性等
  \item \textbf{频率派视角}：理解后验预测检查的校准性质
\end{enumerate}

模型检查不是贝叶斯分析的可选附件，而是其核心组成部分。通过系统地检查模型，我们可以：

\begin{itemize}
  \item 识别模型的不足之处
  \item 指导模型改进方向
  \item 加深对数据和模型的共同理解
  \item 在不确定性量化中保持诚实
\end{itemize}

\section{下节预告}

下一章我们将学习\textbf{模型评估、比较与扩展}（Evaluating, Comparing, and Expanding Models）。在模型检查的基础上，我们将进一步探讨如何系统地比较不同模型的性能，以及如何在必要时扩展或修改模型。

\endinput
